\documentclass[11pt,a4paper]{article}

\usepackage[margin=2.5cm]{geometry}
\usepackage{graphicx}
\usepackage{booktabs}
\usepackage{hyperref}

\title{ICD-10 Category Prediction from Clinical Literals}
\author{Team 10: Phoebe Iglesias, David Redrejo, Pau Rossell}
\date{Fundamentals of Natural Language / NLP-I, 2025--2026}

\begin{document}
\maketitle

\begin{abstract}
This report describes our approach to the UAB-ASHO AI Codification competition.
The task is to predict exactly one ICD category prefix, \texttt{y\_category},
from each short clinical literal.
\end{abstract}

\section{Introduction}
Placeholder for the motivation, competition setting, and why ICD coding matters.

\section{Data and Annotation Analysis}
Placeholder for the first technical phase: analyzing the data and annotations.

\section{Main Challenges}
Placeholder for the second technical phase: understanding the main challenges of
short clinical literals, ambiguity, multilingual variation, and label imbalance.

\section{Methods}
Placeholder for reference methods, baselines, RoBERTa backbone models, and
improved experiments.

\section{Evaluation}
Placeholder for validation metrics, error analysis, interpretability, and
submission provenance.

\section{Conclusion}
Placeholder for final findings and lessons learned.

\bibliographystyle{plain}
\bibliography{references}

\end{document}

